% BHU PYQ -- Professional Exam Template
\documentclass[11pt,a4paper]{article}

\usepackage[a4paper,top=2.5cm,bottom=2.5cm,left=2.8cm,right=2.8cm,headheight=14pt]{geometry}
\usepackage{amsmath,amssymb}
\usepackage[T1]{fontenc}
\usepackage[utf8]{inputenc}
\usepackage{lmodern}
\usepackage{microtype}
\usepackage{fancyhdr}
\usepackage{enumitem}
\usepackage{listings}
\usepackage{hyperref}
\usepackage{lastpage}
\usepackage{parskip}

\hypersetup{colorlinks=true,urlcolor=black,linkcolor=black,
  pdftitle={BVCA-205: Skill Training Python -- BHU PYQ}}

\pagestyle{fancy}
\fancyhf{}
\renewcommand{\headrulewidth}{0.4pt}
\renewcommand{\footrulewidth}{0.4pt}
\fancyhead[L]{\small\textit{Banaras Hindu University}}
\fancyhead[R]{\small\textit{BVCA-205: Skill Training Python}}
\fancyfoot[C]{\small Page \thepage\ of \pageref{LastPage}}

\lstset{
  basicstyle=\small\ttfamily,
  frame=single,
  breaklines=true,
  columns=flexible,
  keepspaces=true
}

\newlist{parts}{enumerate}{2}
\setlist[parts,1]{label=(\alph*),leftmargin=*,itemsep=4pt,topsep=3pt}
\setlist[parts,2]{label=(\roman*),leftmargin=*,itemsep=2pt,topsep=2pt}
\newcommand{\pts}[1]{\hfill\textit{[#1]}}

\begin{document}
\begin{center}
  {\large\bfseries BANARAS HINDU UNIVERSITY}\\[2pt]
  {\normalsize B.Voc. Semester II Examination, 2022-23}\\[6pt]
  \rule{0.8\linewidth}{0.5pt}\\[6pt]
  {\large\bfseries Paper: BVCA-205 --- Skill Training Python}\\[4pt]
  {\normalsize Time: Three hours \hfill Maximum Marks: 70}
\end{center}

\rule{\linewidth}{0.4pt}

\smallskip
\noindent\textit{Instructions: Attempt all questions. Symbols have their usual meanings.}

\bigskip

x
o BVCA - 205 : Skill Training in Python ,
: Answer All Sections.” ,
Candidates are required to write their answers in their own words as far as practicable. |
Section-A
Long Answer type questions to be answered in about 500 words each: 12x2=24 Marks
Attempt any two: ,

\medskip\noindent\textbf{Question 1.} What is a‘programming language? Explain the features of Python programming language. , :

\medskip\noindent\textbf{Question 2.} What.do you understand by iteration? What are various loop structures supported by
Python programming language. ‘

\medskip\noindent\textbf{Question 3.} What is searching? Explain the Binary search algorithm with a suitable program. —
3 4. What are operators? Explain all the operators supported by’ Python programming
language. — . , .
: , , Section- B .
\_ \_ Attempt any six (Answer in about 200 words each): “6x65 36 Marks

\medskip\noindent\textbf{Question 1.} Write a python program to find the sum and average of Nfatifal numbers 3 ee ee
“2. Write a python program to print Fibonacci series upto N. :

\medskip\noindent\textbf{Question 3.} Write a python program to find factorial of a given number using recursion.

\medskip\noindent\textbf{Question 4.} write a python program to check whether a given number is prime or not.
, 5, Write.a python program to create a list with different types of elements and access
: elements of the list using loops. . : .
6. Write a python program to find the largest elementinthe array. . .
7. Write a python program to implement Selection sort. .
‘ 8. Write a python program to perform various operation on strings such as copy, t
concatenate and compare. 7 |
. 9. Write a python program to perform matrix addition. . . . '
‘ : Section- C
Attempt all the questions: Ls ; 10 x 1=10 Marks |
. 1. Which of the following is the correct extension of the Python file?
. a) .python . so
\begin{parts}
  \item pl .
  . d).p .
  a , - .
  @) a
\end{parts}

\medskip\noindent\textbf{Question 2.} What will be the value of the following Python expression? .
4+3\%5
\begin{parts}
  \item 7 : ,
  . b}2 :
  . 4 :
  \item 2 .
  a 3 Which keyword is used for function in Python language? : . '
  \item Function , . ; a |
  \item def . : ae
  : : (}Fun .
  \item Define . . : . .
4. What will be the output of the following Python function?
  . len({"hello",2, 4, 6)) .
  , a) Pfror : . . mo, . ‘
  b}6 ; . . :
  \item 4
d)3 ; co , .
  : 5. How is.a code block indicated in Python? \textasciitilde{} : .
  - a. Brackets.
  \item indentation. .
  \item Key. . :
  \_ o.None of the above. . .
\end{parts}

\medskip\noindent\textbf{Question 6.} What will be the output of the following code snippet? ot \_ ee es
\textasciitilde{} = — example = ["Sunday", “Monday”, "Tuesday", “Wednesday"};
de! example/2] : .
. print(example) ,
\_ a. [‘Sunday’, ‘Monday’, ‘Tuesday’, ‘Wednesday’] — :
\_  b.[ ‘Monday’, ‘Tuesday’, ‘Wednesday’
c. [‘Sunday’, ‘Monday’, ‘Wednesday’] , an . .
\begin{parts}
  \item [‘Sunday’, ‘Tuesday’, ‘Wednesday’] , . ‘ . . :
  \_ 7, What will be the output of the following Python program? . : ‘
  ‘ i=0 .
  while i< 5: : : .
  print(i) ‘ : ,
  i+=1
  ifi==3: . .
  break . .
  . .' else:
  : print(0) . .
  . a) error . .
  b}0120 7
  . c)012 .
  \item none of the mentioned
\end{parts}

\medskip\noindent\textbf{Question 8.} What will be the result of the following expression in Python “2 ** 345 ** 2"? .
\begin{parts}
  \item 65536 ,
  | . c. 169 :
  | d. None of these , :
\end{parts}

\medskip\noindent\textbf{Question 9.} Which of the following statements are used in Exception Handling in Python? - , t
\begin{parts}
  \item try .
  \item except . . :
  \item finally . .
  \item All of the above ,
  “10. Which of the following is a Python tuple? .
  \item , 2,3]
  \item (1, 2, 3) . .
  \item {2, 2, 3}
  . dQ} ; : .
  . ‘ .
  “
  “
  1
\end{parts}

\vfill
\rule{\linewidth}{0.4pt}
\begin{center}\small
\href{https://bhu-exam.vercel.app/}{Click here to visit our website}\\[4pt]\href{https://me-aryan.vercel.app/}{Click here to visit our contributor}\\[4pt]\href{https://quashan-me.vercel.app/}{Click here to visit our contributor}
\end{center}
\end{document}
